\section{Datos}

\subsection{VeRi-776}

El repositorio contiene las particiones \texttt{train}, \texttt{test} y
\texttt{query}, junto con anotaciones XML de color, tipo, identidad y cámara.
Los XML se leen decodificando explícitamente \texttt{gb2312}.

La partición de prueba contiene 200 identidades y se usa por defecto como
fuente de vehículos para evitar fuga de información. La partición de
entrenamiento solo se habilita mediante una opción explícita.

\subsection{Fondos COCO}

Los fondos se leen de COCO val2017 descargado manualmente en
\texttt{data/raw/coco/}. Se excluyen imágenes que contengan automóvil,
motocicleta, autobús, camión o tren. Las restantes se agrupan, por prioridad,
en \texttt{urbano}, \texttt{negativo\_dificil} y \texttt{generico}.

\subsection{Dataset sintético}

El generador usa un lienzo configurable de 640 por 640 píxeles y guarda PNG
sin pérdida, etiquetas YOLO de una sola clase (\texttt{vehicle}) y un
manifiesto JSONL. Selecciona fondos con proporciones configurables, recorta un
cuadrado y lo redimensiona al lienzo.

El tamaño objetivo del vehículo se muestrea con distribución log-uniforme y
se limita al tamaño original para no ampliar los recortes. Cada recorte recibe
una sola operación de remuestreo de alta calidad; el reflejo horizontal se
aplica en la misma etapa de transformación. Por defecto no se aplican filtros
al vehículo, degradación al fondo ni cambios de brillo o contraste.

Las categorías por área final siguen umbrales COCO configurables: pequeño
($<32^2$), mediano ($32^2$--$96^2$) y grande ($>96^2$). El manifiesto registra
la imagen, identidad, atributos heredados, tamaños, escala, posición, reflejo,
grupo de fondo, opciones y semilla.

Limitaciones actuales: no hay máscaras; el rectángulo conserva parte del fondo
original de VeRi y puede producir bordes artificiales. Además, el desajuste de
calidad entre vehículo y fondo puede crear un atajo para el detector. Las
máscaras quedan como punto de extensión futuro.
